\section{Le robot est-il au bord du cadre$\:$?}
Nous avons vu dans le chapitre précédent que les fonctions booléennes {\idt tn ()}, {\idt tne ()}, {\idt te ()}, {\idt tse ()}, {\idt ts ()},
 {\idt tsw ()}, {\idt tw ()}, {\idt tnw ()} peuvent être utilisées dans les instructions
 de test ``if'' pour connaître l'orientation du \robot. Une autre  fonction booléenne,
 {\idt tb ()}, permet de \underline{t}ester si
 le \robot\ est au \underline{b}ord du cadre. Plus précisément {\idt tb ()} vaut {\idt true} (vrai)
 si l'exécution supposée de la  commande {\idt av ()} fait sortir
 le \robot\ de son cadre de déplacement.  L'expression booléenne
 {\idt tb ()} vaut {\idt false} (faux) si l'exécution supposée
 de la  commande {\idt av ()} ne déplace le \robot\ qu'à l'intérieur de son cadre de déplacement.
Par exemple, on peut
 définir une commande {\idt avSansSortir ()} qui ordonne au \robot\ d'exécuter un {\idt av ()} si cela ne le
 fait pas sortir du cadre~:

\input ocaml/AvSansSortirDuCadre.tex

{\sloppy\noindent Les parenthèses de {\mcl not}\ {\idt (tb ())} sont indispensables. Si l'on écrit
{\mcl not}\ {\idt tb ()}, le compilateur comprend que l'on donne deux paramètres {\idt tb}
et {\idt ()} à {\mcl not} et répond `{\tt The function not has type bool -> bool. It is
applied to too many arguments}' (la fonction {\mcl not} reçoit trop de paramètres). Si l'on
oublie les parenthèses vides en écrivant {\mcl not}\ {\idt tb}, il répond
`{\tt The value tb has type unit -> bool but an expression was expected of type bool}'
(la valeur {\idt tb} est une fonction alors que l'on attendait un booléen) et suggère
`{\tt Did you forget to provide () as argument?}' (avez-vous oublié de donner {\idt ()}
en paramètre$\:$?).\par}

%\pagebreak
\section{Boucle ``while''}

 Une boucle ``while'' permet de répéter l'exécution d'une séquence d'ins\-truc\-tions zéro
ou plusieurs fois en testant à chaque itération s'il faut s'arrêter ou continuer. Le mot anglais
``while'' qui se prononce très approximativement `oua$\cdots$i$\cdots$le' signifie `tant que'. La boucle ``while'' permet de
répéter l'exécution d'une séquence d'instructions tant que la valeur d'une expression booléenne
est {\idt true} (vrai). Par exemple, le programme suivant permet de tracer un trait horizontal au milieu
de l'écran en orientant le \robot\ vers l'est et en répétant des {\idt av ()} tant que le \robot\ ne sort
pas du cadre~:

 \input ocaml/TraitHorizontal.tex

Le résultat de l'exécution de ce programme est présenté en réduction dans la figure ci-dessous,
les cadres de déplacement du \robot\ et d'affichage des messages d'erreurs étant en pointillé~:

 \vspace{0.4cm}

 \includevspacePostScript[4.4cm]{TraitHorizontal}{0.95}{0.95}

 \vspace{0.4cm}

\noindent  Une boucle ``while'' a toujours la forme suivante~:

\index{Boucle ``while''}\index{Expression booléenne}
\begin{center}
\begin{tabular}{@{}|p{12.5cm}|@{}} \hline
{\em Boucle``while''~:}\\ \hline \hline
\begin{minipage}[t]{12.5cm}
{\medskip
\mbox{\blc\ \ \ {\mcl while}\hspace{0.3em}\blanc\hspace{0.3em} {\small\rm Expression booléenne} \hspace{0.3em}\blanc\hspace{0.3em}{\mcl do}}\par
\mbox{\blc\ \ \ \ \ \small\rm  \ldots\ Instructions répétées (séparées par des `$;$') \ldots}\par
\mbox{\blc\ \ \ {\mcl done}}\par
\medskip}
\end{minipage}\\ \hline
\end{tabular}
\end{center}

\noindent La  séquence d'instructions séparées par des `$;$' qui est répétée tant que la valeur de
l'{Expression booléenne} vaut {\idt true} (vrai) s'appelle le {\em corps de
la boucle}\index{Corps de boucle}. Comme pour la boucle ``for'', on ne met pas de
point-virgule après la dernière instruction du corps de la boucle, avant le {\mcl done}.
Si l'{Expression booléenne}  vaut {\idt false} quand on commence l'exécution de la
boucle ``while'' alors cette exécution se termine immédiatement sans que le corps de la boucle ne soit
exécuté. Si l'{Expression booléenne}  vaut toujours {\idt true} à chaque itération
alors l'exécution de la boucle ``while'' ne se termine jamais$\:$!
On peut tout de même arrêter l'ordinateur en coupant le courant. Des méthodes plus re\-com\-man\-da\-bles sont expliquées
ci-après.

\underline{Une erreur grave} est d'écrire un corps de boucle qui ne
change pas la valeur de l'Expression booléenne. Par exemple si l'on écrit
{\mcl while}\ {\mcl not}\ {\idt (tb ())}\ {\mcl do}\ {\idt pqtd ()}\ {\mcl done}, le \robot\
tourne sur place sans jamais s'approcher du bord, et l'ordinateur boucle éternellement.

On utilise de préférence les boucles ``for'' (dont l'exécution se termine toujours) en réservant
l'usage des boucles ``while'' au cas où le nombre d'ité\-ra\-tions n'est pas calculable par une expression entière avant de
commencer à exécuter le corps de la boucle.

\begin{exemple}[Cadre]{\rm Le programme ci-dessous permet de tracer un ca\-dre au\-tour d'un dessin~:

 {\raggedright

 \vspace{0.4cm}

 \includevspacePostScript[4.4cm]{TracerCadre}{0.95}{0.95}

 \vspace{0.4cm}

\input ocaml/TracerCadre.tex
 }
\noindent $\Box$ }\end{exemple}

Tous les programmes que nous avions écrits jusqu'alors se terminaient. Ceci veut dire que
 toute exécution du programme prend un temps fini, qui semble parfois long, mais on peut être sûr qu'en attendant suffisamment
 longtemps, on pourra toujours
 observer un résultat. Si par exemple on a levé le crayon du \robot\ et oublié  de le baisser
 par la suite, le résultat est une feuille blanche mais l'exécution du programme se termine tout de
 même.

 Avec la boucle ``while'', on peut écrire des programmes dont l'exécution ne se termine pas.  C'est
 le cas pour le programme suivant~:

 \input ocaml/Boucle1.tex

 \noindent Le \robot\ tourne au centre de l'écran, tant que l'ordinateur marche et que l'on n'intervient pas.
 On peut l'arrêter en tapant, dans la fenêtre du \robot, sur l'une des touches
 \begin{picture}(12,12)(0, 2) %touche d'echappement
  \put(0,0){\framebox(12,12)[lb]{}}
  \put(1.5,2){\makebox(0,0)[lb]{\tiny esc}}
  \end{picture}\hspace{0.3ex},
  \begin{picture}(12,12)(0, 2) %touche §
  \put(0,0){\framebox(12,12)[lb]{}}
  \put(1.5,2){\makebox(0,0)[lb]{\tiny \S}}
  \end{picture}\hspace{0.3ex} ou
  \begin{picture}(12,12)(0, 2) %touche #
  \put(0,0){\framebox(12,12)[lb]{}}
  \put(1.5,2){\makebox(0,0)[lb]{\tiny \#}}
  \end{picture}\hspace{0.3ex}.  L'exécution du programme prend fin car avant
d'exécuter une commande, le \robot\ regarde si l'utilisateur n'a pas tapé sur l'une de
ces touches pour lui demander de s'arrêter.

  Par contre si l'on exécute le programme \ocaml\ suivant~:

 \input ocaml/Boucle2.tex

 \noindent l'exécution boucle sans s'arrêter même si l'on tape sur les touches d'arrêt
  \begin{picture}(12,12)(0, 2) %touche d'echappement
  \put(0,0){\framebox(12,12)[lb]{}}
  \put(1.5,2){\makebox(0,0)[lb]{\tiny esc}}
  \end{picture}\hspace{0.3ex},
  \begin{picture}(12,12)(0, 2) %touche §
  \put(0,0){\framebox(12,12)[lb]{}}
  \put(1.5,2){\makebox(0,0)[lb]{\tiny \S}}
  \end{picture}\hspace{0.3ex} ou
  \begin{picture}(12,12)(0, 2) %touche #
  \put(0,0){\framebox(12,12)[lb]{}}
  \put(1.5,2){\makebox(0,0)[lb]{\tiny \#}}
  \end{picture}\hspace{0.3ex}
puisque le \robot\ ne reçoit
  aucune commande pendant que l'ordinateur exécute cette boucle, dont le corps
  est réduit à {\idt ()}, une instruction qui ne fait rien.

  Pour arrêter l'exécution de ce programme, il faut revenir à la fenêtre de terminal
  dans laquelle on l'a lancé, par exemple avec la commande~:
  \begin{center}{\idt ./robot.sh livre/17\_while/Boucle2.ml}\end{center}
  \noindent maintenir la touche contrôle
\begin{picture}(22,12)(0, 2) %touche Ctrl
\put(0,0){\framebox(22,12)[lb]{}}
\put(1.5,2){\makebox(0,0)[lb]{\tiny Ctrl}}
\end{picture}\hspace{0.3ex}
enfoncée et taper sur la touche
\begin{picture}(12,12)(0, 2) %touche C
\put(0,0){\framebox(12,12)[lb]{}}
\put(1.5,2){\makebox(0,0)[lb]{\tiny C}}
\end{picture}\hspace{0.3ex}. Le système d'exploitation de l'ordinateur interrompt alors
le programme, quoi qu'il soit en train de faire.

Si le programme a été lancé en faisant glisser son icône sur l'application \texttt{LeRobotOCaml.app}, une fenêtre
de terminal s'ouvre, que l'on peut sélectionner avec la souris pour taper Ctrl+C.

%  Sur les premiers Macintosh, pour lesquels ce livre a été écrit à l'origine, il 
%  fallait appuyer sur la touche d'interruption (marquée d'un cercle)
%  placée sur le côté de l'ordinateur (figure \ref{fig:interrupt}) en évitant d'appuyer sur 
%  la touche de réinitialisation (marquée d'un triangle) qui revenait à éteindre puis à 
%  rallumer le Macintosh.
%
%\begin{figure}[htbp]
%  \includevspacePostScript[2.9cm]{Interrupt}{0.95}{0.95}
%  \caption{Touches de réinitialisation et d'interruption des premiers Macintosh}\label{fig:interrupt}
%\end{figure}
  
\clearpage
\noindent \rule{137mm}{0.1mm} %---------------------------------------------
\begin{exercices}{Écrire des programmes \ocaml\  utilisant des boucles ``while'' pour réaliser les
dessins suivants~:

\tiret  la frise de la figure \ref{fig:FriTri} ci-dessous$\:$;
% (page \pageref{fig:FriTri})$\:$;

\tiret  la grille des points d'arrêt du \robot\ de la figure \ref{fig:DgQuGr} (page \pageref{fig:DgQuGr}),
sans utiliser bien entendu la commande {\idt dg ()}$\:$;

\tiret  le quadrillage de la figure \ref{fig:DgQuGr} (page \pageref{fig:DgQuGr})$\:$;

\tiret  le grillage de la figure \ref{fig:DgQuGr} (page \pageref{fig:DgQuGr}) obtenu en superposant deux
qua\-dril\-la\-ges dont un est tourné d'un huitième de tour.

Dans tous les cas le dessin doit être le plus grand possible et s'inscrire à l'intérieur du cadre fixé pour les
déplacements du \robot\ sur l'écran de l'ordinateur.

\noindent $\Box$}\end{exercices}

\vspace*{-1ex}%
\begin{exercice}{Définir les commandes {\idt ag ()},  {\idt ad ()},  {\idt eh ()} et  {\idt eb ()} par des procédures comportant
des boucles ``while'' (on utilisera la fonction booléenne {\idt valBc ()} qui vaut {\idt true} si et
seulement si le crayon du \robot\ est baissé).
\noindent $\Box$}\end{exercice}

\begin{figure}[htbp]
 \includevspacePostScript[4.4cm]{FriseDeTriangles}{0.9}{0.9}
 \caption{Frise de triangles (dans un écran réduit)}\label{fig:FriTri}
\end{figure}
\begin{figure}[htbp]
 \includevspacePostScript[14.6cm]{DgQuadGrillage}{0.95}{0.95}
 \caption{Points d'arrêt, quadrillage et grillage sur un écran réduit}\label{fig:DgQuGr}
\end{figure}

%%%%%%%%%%%%%%%%%%%%%%%%%%%%%%%%%%%%%%%%%%%%%%%%%%%%%%%%%
\vspace*{-1ex}%
\begin{corrige}{
Nous donnons un programme, les autres se trouvant dans le dossier {\idt livre/\allowbreak 17\_while}~:
\input ocaml/FriseDeTriangles.tex
\noindent $\Box$}\end{corrige}
\begin{corrige}{
Nous donnons le corrigé pour la commande {\idt ag}, les autres se trouvant dans le programme
{\idt AgAdEhEb} du dossier {\idt livre/\allowbreak 17\_while}~:
\input ocaml/AgAdEhEb.tex
\noindent $\Box$}\end{corrige}

%%%%%%%%%%%%%%%%%%%%%%%%%%%%%%%%%%%%%%%%%%%%%%%%%%%%%%%%%%%%%%%%%%%%%%%%%%%
\noindent \framebox[\textwidth]{
\begin{minipage}{\myparboxwidth}
\setlength{\parindent}{2em}
\small \sf
\vspace{\abovedisplayskip}
\section*{Addenda}

 Dans d'autres langages de programmation, comme \pascal, on trouve une autre forme de
 répétition non bornée, la boucle ``repeat'' (qui signifie répéter en anglais).
 Elle a la forme suivante~:

 \index{Boucle ``while''}
\begin{center}
\begin{minipage}{10cm}
{\raggedright
\mcl repeat\\
\blc\ \ \small\rm  \ldots\ Instructions répétées (séparées par des `$;$') \ldots\\
\mcl until\hspace{0.3em}\blanc\hspace{0.3em} \small\rm Expression booléenne\\
}\end{minipage}
\end{center}

\noindent Le corps de la boucle est répété ({\idt repeat}) jusqu'à ({\idt until}) ce que la valeur de
 l'{\small\rm Expression booléenne} soit égale à {\idt true}. Il est donc toujours exécuté
 au moins une fois.
%\vspace{\abovedisplayskip}
%\end{minipage}
%}
%%%%%%%%%%%%%%%%%%%%%%%%%%%%%%%%%%%%%%%%%%%%%%%%%%%%%%%%%%%%%%%%%%%%%%%%%%%%

%%%%%%%%%%%%%%%%%%%%%%%%%%%%%%%%%%%%%%%%%%%%%%%%%%%%%%%%%%%%%%%%%%%%%%%%%%%
%\noindent \framebox[\textwidth]{
%\begin{minipage}{\myparboxwidth}
%\setlength{\parindent}{2em}
%\small \sf
%\vspace{\abovedisplayskip}
\indent La boucle ``repeat'' n'existe pas en \ocaml, mais on peut toujours l'exprimer
à l'aide d'une boucle ``while''. Une boucle ``repeat''~:
 {
\parindent \parindentprog
\medskip

\pagebreak[3]
\mbox{\blc  \blk\mcl repeat}\par
\mbox{\blc  \blk\blk\small\rm Corps de la boucle}\par
\mbox{\blc  \blk\mcl until {\small\rm Expression booléenne}}\par

\medskip }

\noindent peut s'écrire en \ocaml\ avec une boucle ``while'' comme suit~:
{
\parindent \parindentprog
\medskip

\mbox{\blc  \blk\small\rm Corps de la boucle$;$}\par
\mbox{\blc  \blk\mcl while\ not\ \idt (\small\rm Expression booléenne\idt )\ \mcl do}\par
\mbox{\blc  \blk\blk\small\rm Corps de la boucle}\par
\mbox{\blc  \blk\mcl done}\par

\medskip }

\noindent  De la même façon, une boucle ``while''~:

{
\parindent \parindentprog
\medskip

\mbox{\blc  \blk\mcl while\ \small\rm Expression booléenne\ \mcl do}\par
\mbox{\blc  \blk\blk\small\rm Corps de la boucle}\par
\mbox{\blc  \blk\mcl done}\par

\medskip }

\noindent pourrait se réécrire avec une boucle ``repeat'' comme suit~:

{
\parindent \parindentprog
\medskip

\mbox{\blc  \blk\mcl if\ \small\rm Expression booléenne\ \mcl then}\par
\mbox{\blc  \blk\blk\mcl repeat}\par
\mbox{\blc  \blk\blk\blk\small\rm Corps de la boucle}\par
\mbox{\blc  \blk\blk\mcl until\ not\ \idt (\small\rm Expression booléenne\idt )}\par

\medskip }

L'usage de l'une ou l'autre forme d'itération ne dépend que du fait que le corps de la boucle
peut ne jamais être exécuté (boucle ``while'') ou doit toujours l'être au moins une fois
(boucle ``repeat''). Les concepteurs d'\ocaml\ ont jugé que la boucle ``while'' suffisait.

\vspace{\abovedisplayskip}
\end{minipage}
}
%%%%%%%%%%%%%%%%%%%%%%%%%%%%%%%%%%%%%%%%%%%%%%%%%%%%%%%%%%%%%%%%%%%%%%%%%%%%
